\documentclass[11pt]{article}
% =============================================================================
%  Task-Relevant Drift: A Causal Account of Catastrophic Forgetting
%  in LoRA Continual Fine-Tuning
%
%  Draft for TMLR (Transactions on Machine Learning Research).
%  All numbers below come from the author's reproducible experiments
%  (observational.csv : N=162, decisive.csv : N=30 paired).
%  Replace the placeholder author block and, if submitting to TMLR, swap the
%  preamble for the official tmlr.sty style file.
% =============================================================================

\usepackage[utf8]{inputenc}
\usepackage[T1]{fontenc}
\usepackage{amsmath,amssymb,amsthm}
\usepackage{graphicx}
\usepackage{booktabs}
\usepackage{microtype}
\usepackage{hyperref}
\usepackage{url}
\usepackage{xcolor}
\usepackage[margin=1in]{geometry}
\usepackage{natbib}

\hypersetup{colorlinks=true, linkcolor=blue, citecolor=blue, urlcolor=blue}

\newcommand{\trdrift}{\mathrm{TR\text{-}drift}}
\newcommand{\totaldrift}{\mathrm{Drift}_{\mathrm{tot}}}
\newcommand{\PA}{\mathcal{P}_A}

\title{\textbf{Task-Relevant Drift: A Causal Account of Catastrophic
Forgetting in LoRA Continual Fine-Tuning}}

\author{
Haithem Barkaoui\\
Recherche indépendante, France\\
\texttt{haythem.barkaoui@gmail.com}
}
\date{\today}

\begin{document}
\maketitle

% -----------------------------------------------------------------------------
\begin{abstract}
Catastrophic forgetting in continual fine-tuning is widely attributed to
\emph{interference} between tasks in a shared parameter substrate, and a
popular family of mitigation methods (gradient orthogonalization, subspace
projection) rests on the assumption that the \emph{static geometric overlap}
between task subspaces predicts how much a model forgets. We test this
assumption directly for low-rank adaptation (LoRA) of GPT-2 and find it
wanting: across 162 continual fine-tuning conditions, static overlap is only a
weak predictor of forgetting ($r=0.39$) and loses its explanatory power once
total parameter motion is controlled for. We instead propose \emph{Task-Relevant
Drift} ($\trdrift$): the component of the post-update parameter change that lies
inside the previous task's functional subspace $\PA$. $\trdrift$ predicts
forgetting substantially better ($r=0.63$) and retains its predictive power
after partialling out total drift. Decomposing total drift into task-relevant
and orthogonal parts further localizes the predictive signal to the
task-relevant component. We then move from correlation to causation with a
controlled decoupling experiment: holding the total backbone step norm fixed,
steering the update \emph{into} $\PA$ rather than orthogonal to it nearly doubles
forgetting (Cohen's $d=0.94$, $p=10^{-4}$), establishing $\trdrift$ as a direct
cause of forgetting rather than a mere correlate. Finally, the relationship
persists across EWC, OGD, and GPM, indicating $\trdrift$ is a general descriptor
rather than an artifact of plain fine-tuning. We release all code and data, and
restrict our claims to GPT-2 small, discussing scaling to larger models as
future work.
\end{abstract}

% -----------------------------------------------------------------------------
\section{Introduction}
\label{sec:intro}

Continual fine-tuning---adapting a pretrained model to a sequence of tasks---is
central to deploying large language models, yet it remains plagued by
\emph{catastrophic forgetting}: performance on earlier tasks degrades as later
tasks are learned \citep{mccloskey1989catastrophic,kirkpatrick2017overcoming}.
Parameter-efficient methods such as LoRA \citep{hu2021lora} reduce the number of
trainable parameters but do not, by themselves, eliminate forgetting, because
the adapted layers still share a common substrate across tasks.

A large body of mitigation methods rests on a single intuitive premise: that
forgetting is driven by \emph{interference} between tasks, and that this
interference is captured by the \emph{static geometric overlap} between the
parameter-space directions each task relies on. Methods that orthogonalize
gradients or project updates into task-specific subspaces are, in effect,
attempts to reduce this static overlap. If the premise holds, static overlap
should predict forgetting.

In this paper we test the premise directly and find that it does not hold in the
form usually assumed. Our contributions are:

\begin{enumerate}
    \item \textbf{A negative result.} Across 162 continual fine-tuning
    conditions, the static overlap between the two tasks' gradient subspaces is
    only weakly correlated with forgetting ($r=0.39$) and is dominated by---and
    largely redundant with---total parameter motion. Static geometry alone does
    not explain forgetting (Section~\ref{sec:negative}).

    \item \textbf{A better predictor.} We define \emph{Task-Relevant Drift}
    ($\trdrift$): the portion of the parameter change that lands inside the
    previous task's functional (gradient-PCA) subspace $\PA$. $\trdrift$
    predicts forgetting more strongly ($r=0.63$) and, unlike static overlap and
    unlike total drift, adds incremental predictive validity: after controlling
    for total drift, partial $r=0.52$ ($p=1.5\times10^{-12}$), with
    $\Delta R^2=+0.21$ (Section~\ref{sec:correlation}).

    \item \textbf{A causal test.} We decouple direction from magnitude: holding
    the total backbone step norm fixed, we steer the update either into $\PA$ or
    orthogonal to it. Motion into $\PA$ causes nearly double the forgetting
    (0.186 vs.\ 0.106; $d=0.94$, $p=10^{-4}$), establishing $\trdrift$ as a
    \emph{cause} of forgetting, not merely a correlate
    (Section~\ref{sec:causal}). This confirms our pre-registered prediction
    that manipulating task-relevant drift independently of total drift changes
    forgetting.
\end{enumerate}

We deliberately scope all claims to GPT-2 small and binary text-classification
tasks; Section~\ref{sec:limitations} discusses why scaling is the key next step.

% -----------------------------------------------------------------------------
\section{Related Work}
\label{sec:related}

\paragraph{Catastrophic forgetting and continual learning.}
Forgetting has been studied for decades \citep{mccloskey1989catastrophic,
french1999catastrophic}. Regularization approaches such as EWC
\citep{kirkpatrick2017overcoming} and LwF \citep{li2017learning} penalize
changes to parameters deemed important for past tasks, while rehearsal/replay
methods \citep{rebuffi2017icarl} retain or regenerate past data. These methods
implicitly or explicitly localize ``important'' directions and try to protect
them.

\paragraph{Subspace and orthogonalization methods.}
A complementary line constrains updates to lie in subspaces that minimize
interference with previous tasks---e.g., orthogonal gradient descent and
gradient projection memory \citep{farajtabar2020orthogonal,saha2021gradient}.
The motivating assumption is that interference is governed by the geometric
overlap between task subspaces. Our negative result questions the predictive
adequacy of \emph{static} overlap and reframes the relevant quantity as the
\emph{realized} motion inside the previous task's functional subspace.

\paragraph{Parameter-efficient fine-tuning.}
LoRA \citep{hu2021lora} and related adapters have become standard. Because the
shared backbone is still perturbed during sequential adaptation, forgetting
persists; our analysis treats the backbone drift as the mechanistic locus.

% -----------------------------------------------------------------------------
\section{Method}
\label{sec:method}

\paragraph{Setup.}
We fine-tune GPT-2 small (124M parameters) for binary sequence classification
with LoRA (rank 16, $\alpha=32$) applied to the attention projections, while
keeping the transformer backbone trainable so that backbone drift---our object
of study---can occur. We train Task~A, then Task~B, and measure how much Task~A
performance degrades. Optimization uses AdamW; the LoRA and backbone learning
rate is swept, the classification head uses a fixed rate. Each task uses 350
training and 150 test examples; we use five seeds.

\paragraph{Forgetting.}
Let $\mathrm{Acc}_A^{(A)}$ be Task~A accuracy immediately after training Task~A,
and $\mathrm{Acc}_A^{(B)}$ its accuracy after subsequently training Task~B. We
define
\begin{equation}
\mathrm{Forgetting} = \max\!\left(0,\; \mathrm{Acc}_A^{(A)} - \mathrm{Acc}_A^{(B)}\right).
\end{equation}

\paragraph{Task functional subspace $\PA$.}
We estimate the directions Task~A relies on from the principal components of its
backbone gradients. We collect backbone gradients over $n$ minibatches of
Task~A, center them, and take the top-$k$ right singular vectors per analyzed
tensor as an orthonormal basis $U_A=\{u_i\}_{i=1}^k$ spanning $\PA$. (In our
runs $n=32$, $k=16$.)

\paragraph{Total drift and Task-Relevant Drift.}
Let $\theta_A$ be the backbone parameters after Task~A and $\theta_B$ after
Task~B, and $\Delta\theta=\theta_B-\theta_A$. We define
\begin{align}
\totaldrift &= \frac{\lVert \Delta\theta \rVert}{\lVert \theta_A \rVert}, &
\trdrift &= \frac{\lVert \Pi_{\PA}\,\Delta\theta \rVert}{\lVert \theta_A \rVert}
= \frac{\sqrt{\sum_i \langle \Delta\theta, u_i\rangle^2}}{\lVert \theta_A \rVert},
\end{align}
where $\Pi_{\PA}$ is the orthogonal projector onto $\PA$. Thus $\trdrift$ is the
share of parameter motion that lands inside the previous task's functional
subspace; computing it requires only inner products with the stored basis.

\paragraph{Static overlap (the tested baseline predictor).}
To operationalize the ``interference = static geometric overlap'' premise, we
also estimate Task~B's functional subspace $\mathcal{P}_B$ (top-$k$ PCs of
Task~B gradients at the post-A point) and measure the overlap between
subspaces as the mean squared cosine between their bases,
\begin{equation}
\mathrm{overlap}(\PA,\mathcal{P}_B) = \frac{1}{k}\,\big\lVert U_A U_B^\top \big\rVert_F^2 \in [0,1],
\end{equation}
which is $1$ for identical subspaces and $0$ for orthogonal ones. This is a
purely \emph{static} quantity, computed before the interfering update.

\paragraph{Causal decoupling.}
To separate direction from magnitude, we re-train Task~B while controlling the
backbone update direction. At each step we project the backbone gradient either
onto $\PA$ (\textsc{into}) or onto its orthogonal complement (\textsc{orth}),
then rescale so that the per-step backbone update has the \emph{same} norm in
both arms. The classification head and LoRA parameters train normally and
identically across arms, so the arms differ \emph{only} in whether backbone
motion is task-relevant, at matched total drift. Any forgetting difference is
therefore attributable to direction relative to $\PA$.

% -----------------------------------------------------------------------------
\section{Experiment 1: Static Overlap Does Not Explain Forgetting}
\label{sec:negative}

We run an observational sweep over five task pairs
(IMDb$\to$AG-News, AG-News$\to$IMDb, SST-2$\to$AG-News, Yelp$\to$AG-News,
IMDb$\to$SST-2), five seeds, and six learning rates
($5\!\times\!10^{-6}$ to $5\!\times\!10^{-4}$), for $N=162$ conditions. For each
condition we record forgetting, static overlap, total drift, and $\trdrift$.

Table~\ref{tab:corr} reports Pearson correlations with forgetting. Static
overlap correlates only weakly ($r=0.39$, 95\% CI $[0.24, 0.52]$)---far below
$\trdrift$ ($r=0.63$)---and, as Section~\ref{sec:correlation} shows, contributes
nothing once magnitude is controlled. The geometric-overlap premise, in the
static form assumed by projection methods, is thus a poor predictor of how much
a model actually forgets.

\begin{figure}[t]
\centering
\includegraphics[width=0.62\linewidth]{fig_overlap_scatter.png}
\caption{Static subspace overlap is a weak, scattered predictor of forgetting
($r=0.39$, $N=162$). Contrast with the tighter task-relevant relationship in
Figure~\ref{fig:trdrift}.}
\label{fig:overlap}
\end{figure}

\begin{table}[t]
\centering
\caption{Correlation of candidate predictors with forgetting ($N=162$).
$\trdrift$ is the strongest. 95\% bootstrap CIs in brackets.}
\label{tab:corr}
\begin{tabular}{lcc}
\toprule
Predictor & Pearson $r$ & 95\% CI \\
\midrule
Static overlap         & $0.39$ & $[0.24,\,0.52]$ \\
Total drift            & $0.47$ & $[0.33,\,0.60]$ \\
Task-Relevant Drift    & $\mathbf{0.63}$ & $[0.48,\,0.74]$ \\
\bottomrule
\end{tabular}
\end{table}

% -----------------------------------------------------------------------------
\section{Experiment 2: Task-Relevant Drift Has Incremental Validity}
\label{sec:correlation}

A predictor is only interesting if it explains variance \emph{beyond} the
obvious confound---here, total parameter motion, which trivially grows with the
learning rate. We therefore test whether $\trdrift$ predicts forgetting after
partialling out total drift.

\begin{figure}[t]
\centering
\includegraphics[width=0.62\linewidth]{fig_trdrift_scatter.png}
\caption{Task-Relevant Drift predicts forgetting ($r=0.63$, $N=162$). Each point
is one continual fine-tuning condition; the relationship is visibly tighter than
for static overlap (Figure~\ref{fig:overlap}).}
\label{fig:trdrift}
\end{figure}

The partial correlation of $\trdrift$ with forgetting, controlling for total
drift, is $r=0.52$ ($p=1.5\times10^{-12}$). A nested regression confirms the
gain: $R^2$ rises from $0.22$ (total drift only) to $0.43$ (total drift
$+\,\trdrift$), an increment of $\Delta R^2=+0.21$, with a positive and highly
significant $\trdrift$ coefficient ($p<10^{-6}$). A mixed-effects model with a
random intercept per source task confirms the effect is not an artifact of
task grouping ($\trdrift$ coefficient $p<10^{-6}$). Static overlap, by contrast,
provides no comparable increment. Table~\ref{tab:incremental} summarizes.

\begin{table}[t]
\centering
\caption{Incremental validity of $\trdrift$ over total drift ($N=162$).}
\label{tab:incremental}
\begin{tabular}{lc}
\toprule
Quantity & Value \\
\midrule
Partial $r$ ($\trdrift$ vs.\ forgetting $\mid$ total drift) & $0.52$ \\
\quad $p$-value & $1.5\times10^{-12}$ \\
$R^2$ (total drift only)              & $0.22$ \\
$R^2$ (total drift $+\,\trdrift$)     & $0.43$ \\
$\Delta R^2$                          & $+0.21$ \\
\bottomrule
\end{tabular}
\end{table}

\paragraph{Robustness to subspace rank $k$.}
TR-drift depends on a chosen rank $k$ for the functional subspace $\PA$. To
verify our finding is not an artifact of a lucky choice of $k$, we re-ran the
observational sweep at $k\in\{8,16,32\}$. The TR-drift--forgetting correlation is
high and significant at every rank: $r=0.65$ ($k=8$), $r=0.70$ ($k=16$), and
$r=0.73$ ($k=32$), all with $p<10^{-7}$ ($N=54$ conditions each). The effect is
thus robust to the choice of subspace rank, and if anything strengthens mildly
as $k$ grows.

% -----------------------------------------------------------------------------
\section{Experiment 2b: Decomposing Drift into Task-Relevant and Orthogonal Components}
\label{sec:decomposition}

The incremental-validity result already hints that not all parameter motion is
equally harmful. We make this precise by decomposing total drift into two
orthogonal parts: the task-relevant component inside $\PA$ ($\trdrift$) and the
remaining orthogonal component
$\mathrm{Drift}_{\perp}=\sqrt{\totaldrift^{2}-\trdrift^{2}}$. We then ask which
component carries the predictive signal for forgetting.

Marginally, both components correlate with forgetting ($\trdrift$: $r=0.63$;
$\mathrm{Drift}_{\perp}$: $r=0.47$), but this is misleading: the orthogonal
component is numerically dominated by---and nearly collinear with---total drift.
The decisive analysis is a joint regression. Regressing forgetting on both
components simultaneously, the task-relevant coefficient is large, positive, and
highly significant, whereas the orthogonal coefficient is small and, after
controlling for $\trdrift$, does \emph{not} carry the same positive association
with forgetting (Table~\ref{tab:decomp}). Equivalently, the partial correlation
of $\trdrift$ with forgetting (controlling for orthogonal drift) is $+0.52$
($p=1.1\times10^{-12}$), while the partial correlation of orthogonal drift
(controlling for $\trdrift$) is $-0.24$ ($p=2.6\times10^{-3}$).

\begin{table}[t]
\centering
\caption{Joint regression of forgetting on the two drift components ($N=162$).
The predictive signal is concentrated in the task-relevant component.}
\label{tab:decomp}
\begin{tabular}{lccc}
\toprule
Component & Coefficient & $p$-value & Partial $r$ \\
\midrule
$\trdrift$ (in $\PA$)               & $+127.1$ & $1.4\times10^{-12}$ & $+0.52$ \\
$\mathrm{Drift}_{\perp}$ (orthogonal) & $-5.6$  & $2.6\times10^{-3}$  & $-0.24$ \\
\midrule
\multicolumn{4}{l}{$R^2 = 0.43$} \\
\bottomrule
\end{tabular}
\end{table}

We are deliberately cautious in interpretation: this analysis does not show that
orthogonal motion is ``safe'' in any absolute sense. What it shows is that,
once the task-relevant component is isolated, the orthogonal component does not
carry the same positive explanatory signal for forgetting. Within the scope of
GPT-2 small and our task suite, the predictive signal is concentrated in motion
inside the previous task's functional subspace---exactly the quantity our causal
experiment (Section~\ref{sec:causal}) manipulates.

% -----------------------------------------------------------------------------
\section{Experiment 3: Task-Relevant Drift Causes Forgetting}
\label{sec:causal}

Correlation, however robust, cannot establish that task-relevant motion
\emph{causes} forgetting rather than co-varying with some third factor. Our
decoupling experiment (Section~\ref{sec:method}) provides the causal test:
holding the total backbone step norm fixed, we steer the update \textsc{into}
$\PA$ versus \textsc{orth} to it, across three strongly-forgetting task pairs,
two learning rates, and five seeds ($N=30$ paired comparisons).

The manipulation succeeds: total drift is matched across arms
(\textsc{into} $=0.0023$, \textsc{orth} $=0.0025$), while by construction the
\textsc{orth} arm has essentially zero $\trdrift$. The effect on forgetting is
large and consistent: the \textsc{into} arm forgets $0.186$ on average versus
$0.106$ for \textsc{orth}---a paired difference of $+0.080$ with 95\% bootstrap
CI $[+0.050, +0.109]$, Cohen's $d=0.94$, Wilcoxon $p=10^{-4}$. The confidence
interval excludes zero, so at matched total motion, steering the backbone
\emph{into} the previous task's functional subspace causes substantially more
forgetting (Table~\ref{tab:causal}).

\begin{figure}[t]
\centering
\includegraphics[width=0.5\linewidth]{fig_causal_intervention.png}
\caption{Causal decoupling at matched total drift ($N=30$). Steering the backbone
\emph{into} $\PA$ produces $0.186$ forgetting versus $0.106$ when steered
orthogonally---nearly double---despite identical total motion ($d=0.94$,
$p=10^{-4}$). Error bars are s.e.m.}
\label{fig:causal}
\end{figure}

\begin{table}[t]
\centering
\caption{Causal decoupling at matched total drift ($N=30$ paired). Motion into
$\PA$ causes more forgetting.}
\label{tab:causal}
\begin{tabular}{lcc}
\toprule
 & \textsc{into} $\PA$ & \textsc{orth} $\PA$ \\
\midrule
Total drift (matched)  & $0.0023$ & $0.0025$ \\
Forgetting             & $\mathbf{0.186}$ & $0.106$ \\
\midrule
\multicolumn{3}{l}{Paired diff $=+0.080$, 95\% CI $[+0.050,+0.109]$,
$d=0.94$, $p=10^{-4}$} \\
\bottomrule
\end{tabular}
\end{table}

% -----------------------------------------------------------------------------
\section{Experiment 4: Task-Relevant Drift Predicts Forgetting Across Mitigation Methods}
\label{sec:baselines}

A natural question is whether $\trdrift$ is merely an artifact of plain
fine-tuning, or a general property that also governs forgetting under
established continual-learning methods. We therefore run three standard
mitigations on the same strongly-forgetting task pairs (five seeds each):
EWC \citep{kirkpatrick2017overcoming}, which penalizes changes to
Fisher-important parameters; OGD \citep{farajtabar2020orthogonal}, which
projects updates orthogonal to stored Task-A gradient directions; and GPM
\citep{saha2021gradient}, which projects orthogonal to the top gradient
subspace. (These are standard reimplementations applied to the backbone target
tensors; we did not exhaustively tune each method's hyperparameters, e.g.\ the
EWC penalty strength.)

Two findings emerge (Table~\ref{tab:baselines}). First, in our two-task,
GPT-2-small setting none of the three methods reduces forgetting by a large
margin relative to vanilla fine-tuning---the differences are small. We do not
over-claim here; reducing forgetting was not the purpose of this experiment.
Second, and more importantly, $\trdrift$ remains correlated with forgetting
\emph{pooled across all four methods} ($r=0.53$, $p=1.3\times10^{-5}$, $N=60$).
The relationship between task-relevant drift and forgetting is thus not specific
to plain fine-tuning: it persists under regularization- and projection-based
mitigations, supporting $\trdrift$ as a general descriptor of forgetting rather
than an artifact of a single training procedure.

\begin{table}[t]
\centering
\caption{Forgetting under mitigation methods ($N=60$; mean $\pm$ s.e.m.).
Differences are small; the key point is that $\trdrift$ tracks forgetting across
\emph{all} methods ($r=0.53$, $p=1.3\times10^{-5}$).}
\label{tab:baselines}
\begin{tabular}{lcc}
\toprule
Method & Forgetting & Mean $\trdrift$ \\
\midrule
Vanilla fine-tuning & $0.262 \pm 0.032$ & $0.0011$ \\
EWC                 & $0.265 \pm 0.025$ & $0.0011$ \\
OGD                 & $0.251 \pm 0.026$ & $0.0007$ \\
GPM                 & $0.240 \pm 0.028$ & $0.0010$ \\
\bottomrule
\end{tabular}
\end{table}

% -----------------------------------------------------------------------------
\section{Discussion}
\label{sec:discussion}

Our three experiments tell a consistent story. The static geometric overlap
between task subspaces---the quantity that motivates many projection-based
mitigations---is a weak and largely redundant predictor of forgetting. What
matters is the \emph{realized} motion inside the previous task's functional
subspace: $\trdrift$ predicts forgetting beyond total drift, and a controlled
manipulation shows it does so causally.

\paragraph{Implications for mitigation.}
The result is encouraging for the \emph{spirit} of orthogonalization methods
(direction matters) while cautioning against relying on \emph{static} overlap
estimates: it is the direction actually traversed during the interfering update,
not the a priori geometric overlap, that drives forgetting. Methods that monitor
or constrain realized task-relevant motion online may be more effective than
those that pre-compute static subspace overlaps.

\paragraph{$\trdrift$ as an online diagnostic.}
Because $\trdrift$ is computed from inner products with a stored low-rank basis,
it can be tracked cheaply during training as an early-warning signal of
impending forgetting, independent of whether one intervenes.

\paragraph{From diagnosis to prevention (a cautionary preliminary).}
Given the causal result, an obvious idea is to \emph{prevent} forgetting by
constraining $\trdrift$ during training---a ``TR-Guard''. We ran a small
preliminary probe (two task pairs, three seeds) comparing vanilla fine-tuning
against a strict static projection that removes the in-$\PA$ component of every
backbone update, and against a soft penalty
$\lambda\lVert\Pi_{\PA}(\Delta\theta)\rVert^2$. Neither variant produced a
statistically reliable reduction in forgetting at this small sample: the strict
projection did not help on average and was highly pair-dependent, while the
penalty showed a promising but non-significant mean reduction
($p\approx0.17$, untuned $\lambda$). We report this honestly as a negative
preliminary: \emph{eliminating static task-relevant drift alone does not
straightforwardly eliminate forgetting}. This does not contradict the causal
finding---moving into $\PA$ does cause forgetting---but indicates that a
\emph{static} one-shot estimate of $\PA$ is insufficient as a control target,
and that effective prevention likely requires a dynamically refreshed subspace
and/or a tuned penalty. We leave a full TR-Guard method to future work.

% -----------------------------------------------------------------------------
\section{Limitations}
\label{sec:limitations}

We restrict all claims to GPT-2 small and to pairs of binary text-classification
tasks; whether the same causal relationship holds at the scale of modern
billion-parameter LLMs is an open and important question, and is the natural
next experiment. Our task sequences are length two; longer continual sequences
may introduce cumulative effects we do not capture. The functional subspace
$\PA$ is estimated from gradient PCA over a finite sample and a chosen rank $k$;
our robustness check over $k\in\{8,16,32\}$ shows the main correlation is stable,
but we did not vary the gradient-sample size. Finally, our forgetting metric is
accuracy-based; loss- or calibration-based notions may behave differently.
None of these caveats undermines the within-scope causal finding, but each marks
a boundary of what we claim.

% -----------------------------------------------------------------------------
\section{Conclusion}
\label{sec:conclusion}

We asked whether the static geometric overlap between task subspaces explains
catastrophic forgetting in LoRA continual fine-tuning, and found that it does
not. The portion of parameter motion that lands inside the previous task's
functional subspace---Task-Relevant Drift---is a stronger predictor, carries
incremental validity beyond total motion, and, in a controlled decoupling
experiment, \emph{causes} more forgetting at matched total drift. Direction in
parameter space, relative to the previous task's functional subspace, is thus a
genuine cause of forgetting. We hope $\trdrift$ proves useful both as a lens on
why continual fine-tuning fails and as a cheap online diagnostic, and we view
scaling the causal test to larger models as the most important next step.

% -----------------------------------------------------------------------------
\section*{Reproducibility Statement}
All experiments are produced by a single script, run end-to-end with one
command. We release the code and the two result files
(\texttt{observational.csv}, $N=162$; \texttt{decisive.csv}, $N=30$; and
\texttt{k\_ablation.csv} for the rank-robustness check) used for
every number and figure in this paper. Random seeds, learning-rate grids, task
pairs, and all hyperparameters are specified in the configuration and logged.
Code and data: \url{https://github.com/Haythemsys/trdt-forgetting}.

% -----------------------------------------------------------------------------
\bibliographystyle{plainnat}
\begin{thebibliography}{99}

\bibitem[Farajtabar et al.(2020)]{farajtabar2020orthogonal}
M.~Farajtabar, N.~Azizan, A.~Mott, and A.~Li.
\newblock Orthogonal gradient descent for continual learning.
\newblock In \emph{AISTATS}, 2020.

\bibitem[French(1999)]{french1999catastrophic}
R.~M. French.
\newblock Catastrophic forgetting in connectionist networks.
\newblock \emph{Trends in Cognitive Sciences}, 3(4):128--135, 1999.

\bibitem[Hu et al.(2021)]{hu2021lora}
E.~J. Hu, Y.~Shen, P.~Wallis, Z.~Allen-Zhu, Y.~Li, S.~Wang, L.~Wang, and
W.~Chen.
\newblock {LoRA}: Low-rank adaptation of large language models.
\newblock \emph{arXiv:2106.09685}, 2021.

\bibitem[Kirkpatrick et al.(2017)]{kirkpatrick2017overcoming}
J.~Kirkpatrick, R.~Pascanu, N.~Rabinowitz, et al.
\newblock Overcoming catastrophic forgetting in neural networks.
\newblock \emph{PNAS}, 114(13):3521--3526, 2017.

\bibitem[Li \& Hoiem(2017)]{li2017learning}
Z.~Li and D.~Hoiem.
\newblock Learning without forgetting.
\newblock \emph{IEEE TPAMI}, 40(12):2935--2947, 2017.

\bibitem[McCloskey \& Cohen(1989)]{mccloskey1989catastrophic}
M.~McCloskey and N.~J. Cohen.
\newblock Catastrophic interference in connectionist networks.
\newblock \emph{Psychology of Learning and Motivation}, 24:109--165, 1989.

\bibitem[Rebuffi et al.(2017)]{rebuffi2017icarl}
S.-A. Rebuffi, A.~Kolesnikov, G.~Sperl, and C.~H. Lampert.
\newblock {iCaRL}: Incremental classifier and representation learning.
\newblock In \emph{CVPR}, 2017.

\bibitem[Saha et al.(2021)]{saha2021gradient}
G.~Saha, I.~Garg, and K.~Roy.
\newblock Gradient projection memory for continual learning.
\newblock In \emph{ICLR}, 2021.

\end{thebibliography}

\end{document}
